\scanpage{233}
\mspage{119\deleted{\unsure{1|8}}}
deux types de sous-ensembles remarquables,
d'une part les \uline{étoiles} (indépendantes de la
donnée de $Z$). d'autre part les \uline{rayons} (ne
dépendant par contre que de la donnée de
$Z_{S}\subset\unsure{\mathfrak S|\mathfrak G}_S$). La relation \inserted{d'intersection} des rayons entre
eux est évidente (deux rayons distinct ont une
intersection vide), alors que celle de deux
étoiles peut-être plus complexe (il faudrait
l'étudier systématiquement \deleted{\unsure{?|$g$}}, \unsure{dans|donc} des cas où les
données structurales sont assez riches peut-être pour
faire une étude \deleted{\unsure{plus}} au moins suivie, p.ex.
le cas associé à un \deleted{\unsure{$E,S$}} $\Lambda$-module libre de
rang 2\ldots). La relation entre étoiles et rayons
est simple, dans le cas justement de l'hyp.
simplificatrice (41) et reformulations (41') (41''), (44).
\deleted{La considérant par} \deleted{Elle signifie en}
Notons \deleted{que ces conditions} \inserted{en effet} \deleted{\unsure{que, $s\in$}}
que si $s\in S.$ $s^\vee\in S^\vee$, et si $E$ \inserted{$\subset S$} est l'étoile définie par $s^\vee$,
alors
\[
 s\in E\quad\Longleftrightarrow\quad(s,s^\vee)\in R
\]
donc \deleted{l'hyp.\unsure{''}}

(48) \deleted{L'hyp.} \inserted{La condition} (41) signifie aussi que l'intersection dans $S$
d'une étoile \inserted{$E$} et d'un rayon \inserted{$D$} est, soit vide,
soit réduite à un point.

Dans le cas où on part d'une structure de graphe,
et où $E$ est l'étoile de $s\in E$, alors $E\cap D\ne\emptyset$
sss les rayons $D$ et \deleted{$\unsure{Z.s}$} $D'=Z.s$
sont incidents.
% archive sheet number: 232

\mspage{119 bis}
\uline{NB} La notion de prégraphe $\Gamma=(S,S^\vee,R)$ avec groupe $Z$ d'homothéties virtuelles
\[
 Z\subset\mathrm{Aut}(\Gamma),\qquad Z\text{ ss-gr. comm.}
\]

